\begin{textsample}{POS Dim 6 – df – Score 10.00 – df/df\_em\_31.txt}  \label{ex:f6_pos_124}
OBJETIVOS E PRINCÍPIOS Os Itinerários \textbf{Formativos} devem possibilitar o aprofundamento e a ampliação das aprendizagens relacionadas às competências gerais da BNCC, além de incorporar, nas estratégias educacionais, valores relacionados à ética, liberdade, democracia, justiça social, pluralidade, solidariedade e \textbf{sustentabilidade}, desenvolvendo habilidades que permitam uma visão de mundo ampla e heterogênea para o respeito e o acolhimento à diversidade.
Ademais, esses itinerários devem fomentar a capacidade de tomar decisões e agir nas mais diversas situações, seja na escola, no trabalho ou na vida.
De acordo com o art. 12 das DCNEM, a partir das áreas do conhecimento e da formação técnica e profissional, os Itinerários \textbf{Formativos} devem ser organizados seguindo o contexto do Distrito Federal e as possibilidades de oferta da instituição educacional, considerando : I - linguagens e suas \textbf{tecnologias} : aprofundamento de conhecimentos \textbf{estruturantes} para aplicação de diferentes linguagens em contextos sociais e de trabalho, estruturando arranjos curriculares que permitam estudos em línguas vernáculas, estrangeiras, clássicas e indígenas, Língua Brasileira de Sinais ( LIBRAS ), das artes, design, linguagens \textbf{digitais}, corporeidade, artes cênicas, roteiros, produções literárias, dentre outros, considerando o contexto local e as possibilidades de oferta pelos sistemas de ensino ; II - matemática e suas \textbf{tecnologias} : aprofundamento de conhecimentos \textbf{estruturantes} para aplicação de diferentes conceitos matemáticos em contextos sociais e de trabalho, estruturando arranjos curriculares que permitam estudos em resolução de \textbf{problemas} e análises complexas, funcionais e não-lineares, análise de dados estatísticos e probabilidade, geometria e topologia, robótica, automação, inteligência artificial, programação, jogos \textbf{digitais}, sistemas dinâmicos, dentre outros, considerando o contexto local e as possibilidades de oferta pelos sistemas de ensino ; III - ciências da natureza e suas \textbf{tecnologias} : aprofundamento de conhecimentos \textbf{estruturantes} para aplicação de diferentes conceitos em contextos sociais e de trabalho, organizando arranjos curriculares que permitam estudos em astronomia, metrologia, física geral, clássica, molecular, quântica e mecânica, instrumentação, ótica, acústica, química dos produtos naturais, análise de fenômenos físicos e químicos, meteorologia e climatologia, microbiologia, imunologia e parasitologia, ecologia, nutrição, zoologia, dentre outros, considerando o contexto local e as possibilidades de oferta pelos sistemas de ensino ; IV - ciências humanas e sociais aplicadas : aprofundamento de conhecimentos \textbf{estruturantes} para aplicação de diferentes conceitos em contextos sociais e de trabalho, estruturando arranjos curriculares que permitam estudos em relações sociais, modelos econômicos, processos políticos, pluralidade cultural, historicidade do universo, do homem e da natureza, dentre outros, considerando o contexto local e as possibilidades de oferta pelos sistemas de ensino ; V - formação técnica e profissional : desenvolvimento de programas educacionais inovadores e atualizados que promovam efetivamente a qualificação profissional dos estudantes para o mundo do trabalho, objetivando sua habilitação profissional tanto para o desenvolvimento de vida e carreira quanto para adaptar -se às novas condições ocupacionais e às exigências do mundo do trabalho contemporâneo e suas contínuas transformações, em condições de competitividade, produtividade e inovação, considerando o contexto local e as possibilidades de oferta pelos sistemas de ensino. ( BRASIL, 2018b ).
A oferta dos Itinerários \textbf{Formativos} deve estar ainda em consonância com as metas e estratégias do Plano Distrital de Educação ( PDE ) e com outras normas e políticas educacionais instituídas no Distrito Federal.
EIXOS \textbf{ESTRUTURANTES} Segundo a Portaria nº 1.432, de 28/12/2018, que estabelece os referenciais para a elaboração dos Itinerários \textbf{Formativos}, organizados a partir de quatro eixos \textbf{estruturantes} : “ Investigação Científica ”, “ Processos Criativos ”, “ \textbf{Mediação} e Intervenção Sociocultural ” e “ \textbf{Empreendedorismo} ”.
Tais eixos integram e integralizam os diferentes arranjos de Itinerários \textbf{Formativos}, criando oportunidades para que os estudantes vivenciem experiências educativas associadas à realidade local, promovendo a sua formação pessoal, profissional e cidadã.
Dessa forma, os Itinerários \textbf{Formativos} devem estar organizados de modo a possibilitar a produção de conhecimentos, o processo criativo, a capacidade de intervenção na realidade local e o \textbf{empreendedorismo}.
Os quatro eixos \textbf{estruturantes} são complementares, e é importante que os Itinerários \textbf{Formativos} incorporem e integrem todos eles, a fim de garantir que os estudantes experimentem diferentes situações de aprendizagem e desenvolvam um conjunto diversificado de habilidades relevantes para sua formação integral.
Assim, os estudantes, no decorrer de seu Ensino Médio, deverão realizar pelo menos um Itinerário \textbf{Formativo} completo, passando, necessariamente, por todos os quatro eixos. \textbf{Figura} 1.
Organização curricular do ensino médio : Eixos \textbf{Estruturantes} dos Itinerários \textbf{Formativos}.
Fonte : SUBEB/DIEM.
Cada eixo \textbf{estruturante} dos Itinerários \textbf{Formativos} traz intencionalidades pedagógicas próprias e integradas aos demais. - Investigação Científica : possibilitar o desenvolvimento da capacidade investigativa e da sistematização do conhecimento, por meio de práticas e produções científicas que permitam o aprofundamento dos conceitos fundantes da ciência, a interpretação e compreensão de fenômenos, assim como a proposição de intervenções que considerem as características locais. - Processos Criativos : desenvolver e expandir a capacidade dos estudantes em propor e realizar projetos inovadores, de forma criativa, e possibilitar o aprofundamento dos conhecimentos sobre as artes, a cultura, as diferentes mídias e as ciências.
Esse eixo deve estimular o pensar e a expressão criativa, por meio da elaboração de soluções inovadoras para uma temática ou um \textbf{problema} identificado, utilizando e integrando diferentes linguagens, tais como obras e espetáculos artísticos e culturais, campanhas e peças de comunicação, programas, \textbf{aplicativos}, jogos, robôs, circuitos, entre outros produtos analógicos e \textbf{digitais}. - \textbf{Mediação} e Intervenção Sociocultural : ampliar a capacidade de os estudantes utilizarem seus conhecimentos adquiridos para atuarem como agentes de mudanças e possibilitar a realização de projetos que contribuam para a construção de uma sociedade mais ética, justa, democrática, inclusiva, solidária e sustentável.
Esse eixo deve, também, estimular a convivência, a \textbf{mediação} de conflitos, a atuação sociocultural e o envolvimento dos estudantes nas questões da vida pública, promovendo o engajamento e a mobilização da comunidade escolar em projetos que promovam transformações. - \textbf{Empreendedorismo} : ampliar a capacidade dos estudantes para unir conhecimentos de diferentes áreas com o fim de empreender projetos e se adaptar aos diferentes contextos do mundo do trabalho, estimulando as habilidades relacionadas ao autoconhecimento e ao protagonismo.
Esse eixo deve, ainda, favorecer o desenvolvimento da autonomia, o foco e a determinação para que os estudantes consigam planejar e conquistar objetivos pessoais ou criar empreendimentos voltados à geração de renda via oferta de produtos e serviços, com ou sem o uso de \textbf{tecnologias}.

% matched lemmas: aplicativo, digital, empreendedorismo, estruturante, figurar, formativo, mediação, problema, sustentabilidade, tecnologia
\end{textsample}
